%------------------------------------------------------------------------------%
\begin{question}
  Determine a posição relativa entre as retas e, caso exista, o ponto de interseção
  \[
    r\colon
    \begin{cases}
      x = 1 +  t  \\
      y = 2 -  t  \\
      z = 3 + 2t
    \end{cases}
  \qquad\qquad
    s\colon
    \begin{cases}
      x = 3 - \tau \\
      y =     \tau \\
      z = 7 -2\tau
    \end{cases}
  \]
\end{question}

%------------------------------------------------------------------------------%
\begin{resolution}{Solução}
  Os vetores diretores são
  \[
    \vec{v} = (1, -1, 2)
  \]
  \[
    \vec{w} = (-1, 1, -2)
  \]

  \pause
  Como
  \[
    \vec{w} = -\vec{v}
  \]
  \pause
  as retas são paralelas
\end{resolution}

%------------------------------------------------------------------------------%
\begin{resolution}{Solução}
  Para verificar se são coincidentes,
  \pause
  comparamos os pontos
  \[
    P = (1, 2, 3)
  \]
  \[
    Q = (3, 0, 7)
  \]
  \pause
  Temos
  \[
    Q - P
    \pause
    = (2, -2, 4)
    \pause
    = 2\vec{v}
  \]
  \pause
  Portanto, o ponto $Q$ pertence à reta $r$

  \pause
  Consequentemente, as duas retas são coincidentes
\end{resolution}

%------------------------------------------------------------------------------%
\endinput
